\item[\textbf{Problem 12}]

In this lab exercise, you will work with a BMP image file named
\textbf{pic\_original.bmp} and explore how encryption affects visual data. Your
task is to encrypt the image using OpenSSL in both ECB (Electronic Code Book)
and CBC (Cipher Block Chaining) modes. For example, you might use commands like
\textbf{openssl enc -aes-256-ecb -in pic\_original.bmp -out pic\_ecb.enc} for
ECB mode and a similar command with \textbf{-aes-256-cbc} for CBC mode.

Because BMP files rely on the first 54 bytes as header information to display
correctly, you must extract the header from the original image and combine it
with the encrypted data: extract the header with \textbf{head -c 54
pic\_original.bmp $>$ header}, extract the encrypted body from byte 55 onward
using \textbf{tail -c +55 pic\_ecb.enc $>$ body}, and then concatenate them via
\textbf{cat header body $>$ pic\_ecb.bmp} (repeat for the CBC-encrypted file).
Finally, open the resulting BMP files using an image viewer and observe any
visible patterns or clues that might reveal aspects of the original picture.
Please submit screenshots of your terminal commands, the resulting images, and a
brief report discussing your observations and any differences between the ECB
and CBC encryption outcomes.

For Linux and macOS, you can directly use the head and tail commands. For
Windows users, you have two options: either run these commands via \textbf{WSL}
(Windows Subsystem for Linux) or use a hex editor (such as HxD) to manually copy
the first 54 bytes from \textbf{pic\_original.bmp} and overwrite the first 54
bytes in each encrypted file. Then, save the modified files as
\textbf{pic\_ecb.bmp} and \textbf{pic\_cbc.bmp}.

\Solution

The following commands were executed to perform the encryption. The 256-bit
encryption key used for this process was:

{\centering \code{Kr6XO5YCMXlLEZnmPpfK5KE0pyA2MO0Uw2CA1J3YqYA7LLMs9O}}:

\begin{minipage}{\linewidth}
    \centering
    \includegraphics[width=0.8\linewidth]{problem-12-terminal-commands.png}
    \captionof{figure}{Execution of OpenSSL encryption commands}
    \label{problem-12:terminal}
\end{minipage}

Here is the original, unencrypted image for reference:

\begin{minipage}{\linewidth}
    \centering
    \includegraphics[width=0.8\linewidth]{problem-12-pic-original.jpg}
    \captionof{figure}{Original unencrypted image}
    \label{problem-12:pic-original}
\end{minipage}

Here is the resulting image after encryption using Electronic Codebook (ECB)
mode:

\begin{minipage}{\linewidth}
    \centering
    \includegraphics[width=0.8\linewidth]{problem-12-pic-ecb.jpg}
    \captionof{figure}{Image encrypted using ECB mode}
    \label{problem-12:pic-ecb}
\end{minipage}

Here is the resulting image after encryption using Cipher Block Chaining (CBC)
mode:

\begin{minipage}{\linewidth}
    \centering
    \includegraphics[width=0.8\linewidth]{problem-12-pic-cbc.jpg}
    \captionof{figure}{Image encrypted using CBC mode}
    \label{problem-12:pic-cbc}
\end{minipage}

\vspace{1em}
\textbf{Analysis:}

As observed in Figure \ref{problem-12:pic-ecb}, encrypting the image using ECB
mode fails to mask the visual patterns of the original file. This occurs because
ECB mode is deterministic; it encrypts identical plaintext blocks into identical
ciphertext blocks. Consequently, regions of uniform color in the original image
translate directly into repeating ciphertext blocks, leaking the underlying
structural data.

On the other hand, the image encrypted using CBC mode (Figure
\ref{problem-12:pic-cbc}) successfully hides all distinct visual
characteristics. By utilizing an Initialization Vector (IV) and XORing each
plaintext block with the previous ciphertext block before encryption, CBC
diffuses the data. This eliminates block repetition, resulting in a ciphertext
output that appears as high-entropy, pseudorandom noise.
